\chapter{Retas e polígonos}\label{ch:g4:geometry}

Com régua, esquadro e compasso, a geometria vira um ofício. Este
capítulo ensina as duas posições especiais de retas ---
\omterm{def:g4:geometry:def}{perpendiculares} e \omterm{def:g4:geometry:def}{paralelas} --- no nível das ferramentas, e faz um
passeio pela família dos \omterm{def:g4:geometry:polygon}{polígonos}. (A teoria dessas noções vem no
\cref{ch:g6:lines}.)

\section{Perpendiculares e paralelas}

\begin{definition}[Perpendiculares, paralelas]\label{def:g4:geometry:def}
Duas retas são \emph{perpendiculares}\index{retas perpendiculares}
quando se cruzam formando um \omterm{def:g3:shapes:rightangle}{ângulo reto} (conferido com o esquadro,
marcado com um quadradinho). Duas retas são
\emph{paralelas}\index{retas paralelas} quando seguem a mesma direção
e nunca se cruzam --- como os dois trilhos de uma ferrovia, sempre à
mesma distância uma da outra.
\end{definition}

\begin{method}[Traçar com o esquadro]\label{met:g4:geometry:draw}
\emph{Perpendicular a uma reta $d$ passando por um ponto $P$}: apoie
uma das bordas do \omterm{def:g3:shapes:rightangle}{ângulo reto} sobre $d$, deslize até a outra borda
alcançar $P$ e trace.

\emph{Paralela a $d$ passando por $P$}: apoie a borda do esquadro
sobre $d$, encoste a régua em outra borda para servir de trilho,
deslize o esquadro ao longo da régua até $P$ e trace.
\end{method}

\begin{omfigure}
\begin{tikzpicture}[scale=0.85]
  \draw[thick] (-0.4,0) -- (4.2,0) node[below, font=\small] {$d$};
  \draw[omDef, thick] (1.8,-0.9) -- (1.8,2.3);
  \draw[thin] (1.8,0) ++(-0.25,0) -- ++(0,0.25) -- ++(0.25,0);
  \fill (1.8,1.6) circle (1.6pt) node[right, font=\small] {$P$};
  \node[below, font=\small] at (1.9,-1.1) {perpendicular por $P$};
  \begin{scope}[xshift=7.4cm]
  \draw[thick] (-0.4,0) -- (4.2,0) node[below, font=\small] {$d$};
  \draw[omThm, thick] (-0.4,1.6) -- (4.2,1.6);
  \fill (1.6,1.6) circle (1.6pt) node[above, font=\small] {$P$};
  \draw[Stealth-Stealth, thin] (3.4,0) -- (3.4,1.6)
    node[midway, right, font=\small] {mesma distância em toda parte};
  \node[below, font=\small] at (1.9,-1.1) {paralela por $P$};
  \end{scope}
\end{tikzpicture}

{\small As duas construções. \omterm{def:g4:geometry:def}{Retas paralelas} mantêm distância
constante: medir essa distância em dois lugares é uma boa
conferência.}
\end{omfigure}

\begin{example}[Reconhecê-las nas figuras]\label{ex:g4:geometry:spot}
Num retângulo, cada lado é perpendicular aos dois vizinhos e paralelo
ao lado oposto. O papel quadriculado é uma rede pronta de \omterm{def:g4:geometry:def}{paralelas} e
\omterm{def:g4:geometry:def}{perpendiculares} --- é o que o torna tão confortável para desenhar.
\end{example}

\section{Polígonos}

\begin{definition}[Polígono]\label{def:g4:geometry:polygon}
Um \emph{polígono}\index{polígono} é uma figura fechada cujo contorno
é feito de segmentos de reta, os seus \emph{lados}; os pontos dos
cantos são os \emph{vértices}\index{vértice}. Os polígonos recebem nome pelo número
de lados: \emph{triângulo} ($3$), \emph{quadrilátero}\index{quadrilátero} ($4$),
\emph{pentágono}\index{pentágono} ($5$), \emph{hexágono}\index{hexágono} ($6$),
\emph{octógono}\index{octógono}
($8$).
\end{definition}

\begin{omfigure}
\begin{tikzpicture}[scale=0.62]
  \draw[thick, fill=omDef!10] (0,0) -- (2,0) -- (1.2,1.7) -- cycle;
  \node[below, font=\small] at (1,-0.3) {triângulo};
  \begin{scope}[xshift=3.4cm]
  \draw[thick, fill=omThm!10] (0,0) -- (2.2,0) -- (2.5,1.5) --
    (0.5,1.8) -- cycle;
  \node[below, font=\small] at (1.25,-0.3) {quadrilátero};
  \end{scope}
  \begin{scope}[xshift=7.4cm]
  \draw[thick, fill=omMeth!10] (90:1.1) -- (162:1.1) -- (234:1.1) --
    (306:1.1) -- (18:1.1) -- cycle;
  \node[below, font=\small] at (0,-1.4) {pentágono};
  \end{scope}
  \begin{scope}[xshift=10.9cm]
  \draw[thick, fill=omProp!10] (0:1.1) -- (60:1.1) -- (120:1.1) --
    (180:1.1) -- (240:1.1) -- (300:1.1) -- cycle;
  \node[below, font=\small] at (0,-1.4) {hexágono};
  \end{scope}
\end{tikzpicture}

{\small A família dos \omterm{def:g4:geometry:polygon}{polígonos} (o \omterm{def:g4:geometry:polygon}{pentágono} e o \omterm{def:g4:geometry:polygon}{hexágono} desenhados
aqui são \emph{regulares}: lados iguais e cantos iguais). O círculo
não é um \omterm{def:g4:geometry:polygon}{polígono}: o contorno dele não é feito de segmentos.}
\end{omfigure}

\begin{example}[Descrever um polígono]\label{ex:g4:geometry:describe}
Uma brincadeira útil: descrever uma figura com tanta precisão que um
colega consiga redesenhá-la sem vê-la. ``Um \omterm{def:g4:geometry:polygon}{quadrilátero} $ABCD$ com
$AB = 5$ cm, o ângulo em $A$ reto, $AD = 3$ cm\,\dots'' --- é o
vocabulário que faz a mensagem funcionar.
\end{example}

\section{O círculo e o compasso}

\begin{definition}[Círculo]\label{def:g4:geometry:circle}
O \emph{círculo} de centro $O$ e \omterm{def:g3:shapes:shapes}{raio} $3$ cm é o conjunto de todos os
pontos que estão a exatamente $3$ cm de $O$. O compasso o traça:
ponta seca em $O$, abertura de $3$ cm, gire. Um
\emph{diâmetro}\index{diâmetro} é um segmento que passa pelo centro e
liga dois pontos do círculo; ele mede o dobro do \omterm{def:g3:shapes:shapes}{raio}.
\end{definition}

\begin{example}[O compasso como guardador de distâncias]\label{ex:g4:geometry:compass}
O compasso faz mais do que traçar círculos: ele \emph{transporta
distâncias}. Para marcar todos os pontos a $4$ cm de um ponto $A$,
trace o círculo de centro $A$ e \omterm{def:g3:shapes:shapes}{raio} $4$ cm; para comparar dois
segmentos sem régua, abra o compasso sobre um e teste-o no outro.
\end{example}

\begin{method}[Reproduzir uma figura na malha]\label{met:g4:geometry:reproduce}
\begin{enumerate}
  \item Escolha um \omterm{def:g4:geometry:polygon}{vértice} de partida e a posição dele na malha
        vazia;
  \item percorra a figura \omterm{def:g4:geometry:polygon}{vértice} por \omterm{def:g4:geometry:polygon}{vértice}, contando os passos da
        malha (``$3$ para a direita, $2$ para cima'');
  \item ligue os \omterm{def:g4:geometry:polygon}{vértices} em ordem e compare com o modelo: mesmos
        comprimentos, mesmas direções.
\end{enumerate}
\end{method}

\section{Exercícios}

\begin{exercise}[$\star$]\label{exo:g4:geometry:1}
Trace uma reta $d$ e um ponto $P$ fora dela. Construa a perpendicular
a $d$ por $P$ e depois a paralela a $d$ por $P$ (com o método da
régua-trilho). Marque o \omterm{def:g3:shapes:rightangle}{ângulo reto}.
\end{exercise}

\begin{exercise}[$\star$]\label{exo:g4:geometry:2}
Observe as letras maiúsculas E, H, N, T, Z. Em cada uma, encontre um
par de traços paralelos e um par de traços \omterm{def:g4:geometry:def}{perpendiculares}, quando
existirem.
\end{exercise}

\begin{exercise}[$\star$]\label{exo:g4:geometry:3}
Dê o nome de cada \omterm{def:g4:geometry:polygon}{polígono}: $6$ lados; $4$ lados; $3$ lados; $8$
lados; $5$ lados. Quantos \emph{\omterm{def:g4:geometry:polygon}{vértices}} cada um tem?
\end{exercise}

\begin{exercise}[$\star$]\label{exo:g4:geometry:4}
Trace um \omterm{def:g4:geometry:polygon}{pentágono} qualquer (não necessariamente regular) e nomeie os
\omterm{def:g4:geometry:polygon}{vértices} $A$, $B$, $C$, $D$, $E$. Quantos lados ele tem? Trace as
diagonais que saem de $A$ (segmentos que ligam $A$ a \omterm{def:g4:geometry:polygon}{vértices} não
vizinhos): quantas são?
\end{exercise}

\begin{exercise}[$\star$]\label{exo:g4:geometry:5}
Trace um círculo de centro $O$ e \omterm{def:g3:shapes:shapes}{raio} $4$ cm; um \omterm{def:g4:geometry:circle}{diâmetro} $[AB]$; e
um ponto $M$ do círculo fora de $[AB]$. Dê os comprimentos $OA$, $OM$
e $AB$ sem medir.
\end{exercise}

\begin{exercise}[$\star$]\label{exo:g4:geometry:6}
Com o compasso como guardador de distâncias
(\cref{ex:g4:geometry:compass}), compare os lados de um triângulo
traçado por um colega e decida se ele tem dois lados iguais.
\end{exercise}

\begin{exercise}[$\star$]\label{exo:g4:geometry:7}
No papel quadriculado, reproduza este caminho e feche-o num \omterm{def:g4:geometry:polygon}{polígono}:
comece num ponto $A$; ande $4$ para a direita; $2$ para cima; $3$
para a esquerda; $2$ para cima; depois volte a $A$ em linha reta.
Qual é o nome do \omterm{def:g4:geometry:polygon}{polígono} que você obtém (conte os lados)?
\end{exercise}

\begin{exercise}[$\star$]\label{exo:g4:geometry:8}
Verdadeiro ou falso? Explique com um desenho.
``Duas \omterm{def:g4:geometry:def}{retas perpendiculares} sempre se cruzam.''
``Duas \omterm{def:g4:geometry:def}{retas paralelas} nunca se cruzam.''
``Uma reta pode ser paralela a uma reta e perpendicular a outra.''
\end{exercise}

\begin{exercise}[$\star\star$]\label{exo:g4:geometry:9}
Trace uma reta $d$, depois uma reta $e$ perpendicular a $d$, depois
uma reta $f$ perpendicular a $e$. Olhe para $d$ e $f$: o que você
percebe? (Essa observação vira um teorema no \cref{ch:g6:lines}.)
\end{exercise}

\begin{exercise}[$\star\star$]\label{exo:g4:geometry:10}
Descreva a figura abaixo com palavras para um colega, sem mostrá-la:
um quadrado de lado $4$ cm, com um círculo de \omterm{def:g3:shapes:shapes}{raio} $2$ cm centrado no
meio do quadrado (tocando os quatro lados). Escreva a descrição como
instruções numeradas.
\end{exercise}

\begin{exercise}[$\star\star$]\label{exo:g4:geometry:11}
Quantas diagonais tem um \omterm{def:g4:geometry:polygon}{quadrilátero}? E um \omterm{def:g4:geometry:polygon}{pentágono}? E um
\omterm{def:g4:geometry:polygon}{hexágono}? Desenhe os três casos, conte com cuidado e descreva o
padrão que você observa.
\end{exercise}
